\section{Deriving New Cycle Integrals}

These properties describe how to build a new cycle integral from an
existing one: replicating its backing list, shifting its index, rotating
its gaps, filtering out multiples of a value, and reconstructing the
filtered result directly. Properties 6.1, 6.4--6.10 are fully
Stainless-verified; Properties 6.2 and 6.3 have mathematical proofs but are
not yet Stainless-verified.

\begin{itemize}
  \item x-fold cycle expansion: the physical period changes while the
    represented stream is preserved --- Subsection~6.1
  \item Index shifts: right and left --- Subsections~6.2--6.3
  \item Gap rotation with head adjustment: rotating the gap cycle shifts
    the represented integral by one position --- Subsection~6.4
  \item Survivor filtering: exactness and structure of the scan that
    retains non-multiples --- Subsections~6.5--6.6
  \item Filter-merge reconstruction: the resulting gap cycle after removing
    a multiple --- Subsections~6.7--6.10
\end{itemize}

\subsection{x-fold Cycle Expansion}

Let $L^{(x)}$ be the $x$-fold concatenation of a list $L \in \mathbb{L}$.
This operation changes the physical backing cycle, but it does not change
the unbounded stream represented by the cycle.

The values that change are the finite storage properties:

\begin{equation*}
\begin{aligned}
|L^{(x)}| &= x \cdot |L|
  \quad &&\text{[Expanded physical period]} \\
\sum L^{(x)} &= x \cdot \sum L
  \quad &&\text{[Expanded physical sum]}
\end{aligned}
\end{equation*}

{\raggedright
The period equation is verified in
\href{https://github.com/thiagomata/prime-numbers/blob/master/src/main/scala/v1/chapter4/cycle/integral/recursive/properties/RepeatedGapIntegralProperties.scala}{\texttt{RepeatedGapIntegralProperties::\allowbreak assertRepeatedPeriodIsMultiplied}}.
\par}

The values that do not change are the semantic stream properties:

\begin{equation*}
\begin{aligned}
L^{(x)}_i &= L_{(i \operatorname{mod} |L|)}
  \quad &&\text{[Same cycle lookup]} \\
\operatorname{CycleIntegral}(L^{(x)}, init)_i
  &= \operatorname{CycleIntegral}(L, init)_i
  \quad &&\text{[Same integral stream]}
\end{aligned}
\end{equation*}

{\raggedright
The cycle lookup equation is verified in
\href{https://github.com/thiagomata/prime-numbers/blob/master/src/main/scala/v1/chapter4/cycle/integral/recursive/properties/RepeatedGapIntegralProperties.scala}{\texttt{RepeatedGapIntegralProperties::\allowbreak assertReplicatedCycleValueEqual}}.
\par}

So the expansion is not an invariance of the finite cycle object itself. It
is an invariance of the infinite stream that the cycle object represents.

\begin{equation*}
\begin{aligned}
n &= |L|, \quad L = [v_0, \dots, v_{n-1}]
  &&\text{[Length of original cycle]} \\
L^{(x)}
&:= \underbrace{L \mathbin{\texttt{++}} \dots \mathbin{\texttt{++}} L}_{x \text{ copies}}
  &&\text{[Concatenate $x$ copies]} \\
m &:= |L^{(x)}| = x \cdot n
  &&\text{[Length of new cycle]} \\
T &:= \sum_{j=0}^{n-1} v_j
  &&\text{[Original cycle sum]}
\end{aligned}
\end{equation*}

The list-level repetition is defined by applying the original list at the
original index modulo $n$:

\begin{equation*}
\begin{aligned}
L^{(x)}_i &= L_{(i \operatorname{mod} n)}
  \quad &&\text{for } 0 \le i < x \cdot n \\
\sum L^{(x)} &= x \cdot \sum L
  \quad &&\text{[Repeated sum]}
\end{aligned}
\end{equation*}

{\raggedright
These list-level properties are verified by \texttt{RepeatedList} and
\texttt{ListRepeatProperties}: the repeated list size is $x \cdot n$, the
repeated value at index $i$ is the original value at
$i \operatorname{mod} n$, and the repeated sum is $x$ times the original
sum. See
\href{https://github.com/thiagomata/prime-numbers/blob/master/src/main/scala/v1/chapter3/list/RepeatedList.scala}{\texttt{RepeatedList}},
\href{https://github.com/thiagomata/prime-numbers/blob/master/src/main/scala/v1/chapter3/list/properties/RepeatedListProperties.scala}{\texttt{RepeatedListProperties::\allowbreak assertSumMultiplier}},
and
\href{https://github.com/thiagomata/prime-numbers/blob/master/src/main/scala/v1/chapter3/list/properties/ListRepeatProperties.scala}{\texttt{ListRepeatProperties::\allowbreak assertRepeatedIndex}}.
\par}

\begin{equation*}
\begin{aligned}
\operatorname{CycleIntegral}(L^{(x)}, init)_i
&= (i \operatorname{div} m)\cdot T^{(x)} + I^{(x)}_{i \bmod m} + init
  &&\text{[By Definition]} \\
&= (i \operatorname{div} (x \cdot n))\cdot (x \cdot T) + I^{(x)}_{i \bmod (x \cdot n)} + init
  &&\text{[Substitution]} \\
&= (i \operatorname{div} n)\cdot T + I_{i \bmod n} + init
  &&\text{[Exact simplification]} \\
&= \operatorname{CycleIntegral}(L, init)_i
  &&\text{[Exact value reproduction]}
\end{aligned}
\end{equation*}

\begin{equation*}
\therefore \ \forall \ i \in \mathbb{N}_0: \ \operatorname{CycleIntegral}(L^{(x)}, init)_i = \operatorname{CycleIntegral}(L, init)_i \quad \blacksquare
\end{equation*}

{\raggedright
This property is verified in the
\href{https://github.com/thiagomata/prime-numbers/blob/master/src/main/scala/v1/chapter4/cycle/integral/recursive/properties/CycleIntegralProperties.scala}{\texttt{CycleIntegralProperties::\allowbreak assertRepeatedValuesIntegralMatches}}.
The full Scala verification code is in Appendix~A.8.
\par}

\subsection{Right Index Shift}

Let $L' \in \mathbb{L}$ be the right shift of $L \in \mathbb{L}$ by one
position, and $init' := init + L_0$ be the shifted initial value. Then the
CycleIntegral of $L'$ with $init'$ reproduces the CycleIntegral of $L$ with
$init$, shifted by one position.

\begin{equation*}
\begin{aligned}
n &= |L|, \quad L = [v_0, \dots, v_{n-1}] \\
L' &:= [v_1, v_2, \dots, v_{n-1}, v_0] \\
S' &= S \quad &\text{[Cycle Sum Invariance]} \\
init' &:= init + L_0
\end{aligned}
\end{equation*}

\begin{equation*}
\operatorname{CycleIntegral}(L, init)_{i+1} = \operatorname{CycleIntegral}(L', init')_{i} \quad \forall i \in \mathbb{N}_0
\end{equation*}

\textbf{Base Case}:

\begin{equation*}
\begin{aligned}
A &:= \operatorname{CycleIntegral}(L, init)_i \\
B &:= \operatorname{CycleIntegral}(L', init')_{i} \\
A_0 &= init + L_0 \\
A_1 &= init + L_0 + L_1 \\
B_0 &= init' + L'_0 = (init + L_0) + L'_0 = init + L_0 + L_1 = A_1
\end{aligned}
\end{equation*}

\textbf{Induction Step}:

\begin{equation*}
\begin{aligned}
B_{i-1} &= A_i \quad &\text{[Induction Hypothesis]} \\
A_{i+1} &= A_i + L_{((i + 1) \operatorname{mod} n)} \quad &\text{[By Definition]} \\
B_i &= B_{i-1} + L'_{(i \operatorname{mod} n)} \quad &\text{[By Definition]} \\
    &= A_i + L_{((i + 1) \operatorname{mod} n)} \quad &\text{[By Induction Hypothesis]} \\
    &= A_{i+1} \quad &\text{[By Definition]}
\end{aligned}
\end{equation*}

\begin{equation*}
\therefore \ \forall \ i \in \mathbb{N}_0: \ \operatorname{CycleIntegral}(L, init)_{i+1} = \operatorname{CycleIntegral}(L', init')_{i} \quad \blacksquare
\end{equation*}

The one-period \texttt{CycleIntegral} wrapper is verified directly. If the
shifted cycle uses the one-step rotation of the original backing values and
the shifted initial value is advanced by the original first gap, then the
shifted integral at position \texttt{i} equals the original integral at
position \texttt{i + 1} for every stored-period index with
\texttt{i + 1 < period}.

{\raggedright
This property is verified in
\href{https://github.com/thiagomata/prime-numbers/blob/master/src/main/scala/v1/chapter4/cycle/integral/recursive/properties/GapProperties.scala}{\texttt{GapProperties::\allowbreak assertRotateOneCycleIntegralShiftsByOne}}.
The full Scala verification code is in Appendix~A.9.
\par}

The article's mathematical statement above is the all-position version. The
verified lemma proves the stored-period core; packaging the universal
all-position wrapper only needs the already verified full-cycle shift law.

\subsection{Left Index Shift}

Let $L'' \in \mathbb{L}$ be the left shift of $L \in \mathbb{L}$ by one
position ($|L| > 1$), and $init'' := init + L_0 - L_{n-1}$ be the shifted
initial value. Then the CycleIntegral of $L''$ with $init''$ reproduces the
CycleIntegral of $L$ with $init$, shifted by one position in the opposite
direction.

\begin{equation*}
\begin{aligned}
n &= |L|, \quad L = [v_0, \dots, v_{n-1}], \quad n > 1 \\
L'' &:= [v_{n-1}, v_0, v_1, \dots, v_{n-2}] \\
S'' &= S \quad &\text{[Cycle Sum Invariance]} \\
init'' &:= init + L_0 - L_{n-1}
\end{aligned}
\end{equation*}

\begin{equation*}
\operatorname{CycleIntegral}(L, init)_i = \operatorname{CycleIntegral}(L'', init'')_{i+1} \quad \forall i \in \mathbb{N}_0
\end{equation*}

\textbf{Base Case}:

\begin{equation*}
\begin{aligned}
C &:= \operatorname{CycleIntegral}(L'', init'')_{i} \\
C_1 &= init'' + L''_0 = (init + L_0 - L_{n-1}) + L''_0 \\
    &= init + L_0 - L_{n-1} + L_{n-1} = init + L_0 = A_0
\end{aligned}
\end{equation*}

\textbf{Induction Step}:

\begin{equation*}
\begin{aligned}
C_{i+1} &= A_i \quad &\text{[Induction Hypothesis]} \\
A_{i+1} &= A_i + L_{((i + 1) \operatorname{mod} n)} \quad &\text{[By Definition]} \\
C_{i+1} &= C_{i} + L''_{(i \operatorname{mod} n)} \quad &\text{[By Definition]} \\
       &= A_i + L_{((i - 1 + n) \operatorname{mod} n)} \quad &\text{[By Induction Hypothesis]} \\
       &= A_{i+1} \quad &\text{[By Definition]}
\end{aligned}
\end{equation*}

\begin{equation*}
\therefore \ \forall \ i \in \mathbb{N}_0: \ \operatorname{CycleIntegral}(L, init)_{i+1} = \operatorname{CycleIntegral}(L'', init'')_{i} \quad \blacksquare
\end{equation*}

\subsection{Gap Rotation with Head Adjustment}

Rotating a gap cycle by one position and adjusting the head shifts the
entire integral by one position.

\begin{equation*}
\begin{aligned}
&\operatorname{GapList}(\text{head} + \text{gaps}_0,\; \operatorname{tail}(\text{gaps}) \mathbin{\texttt{++}} (\text{gaps}_0 :: L_e))_i \\
  &\quad = \operatorname{GapList}(\text{head},\; \text{gaps})_{i + 1}
  \quad \text{for } 0 \leq i \text{ and } i + 1 < \operatorname{period}(\operatorname{ci}).
\end{aligned}
\end{equation*}

\textbf{Proof.} At $i=0$, the adjusted head is
$\text{head}+\text{gaps}_0$, which is the original integral at position
$1$. For $i>0$, assume the shifted integral at $i-1$ equals the original
integral at $i$. The rotated gap at $i$ is the original gap at $i+1$;
applying the one-step property to both integrals gives

\begin{equation*}
\begin{aligned}
I'_i &= I'_{i-1} + \text{gaps}'_{i-1} \\
     &= I_i + \text{gaps}_i \\
     &= I_{i+1}. \\
\therefore\ I'_i &= I_{i+1}
\quad \blacksquare\ \text{[Q.E.D.]}
\end{aligned}
\end{equation*}

{\raggedright
This property is verified in
\href{https://github.com/thiagomata/prime-numbers/blob/master/src/main/scala/v1/chapter4/cycle/integral/recursive/properties/GapProperties.scala}{\texttt{GapProperties::\allowbreak assertRotateOneShiftsIntegralByOne}}.
It delegates to the verified
\texttt{ShiftedList.assertShiftedApplyIsOriginalPlusOne}; the full source is
linked there rather than repeated inline.
\par}

For a filter value $f$ and a range $[start, start + count)$, the survivor
sequence is the sub-sequence of $ci$'s values at those positions whose
remainder modulo $f$ is nonzero --- the cycle-integral analogue of one
Eratosthenes sieve step~\cite{hardywright1979}, where positions divisible by
$f$ are removed and only the non-multiples remain:

\begin{equation*}
\begin{aligned}
&S := \operatorname{survivorValues}(ci, f, start, count) := \\
&\quad \big[\, ci(pos) \mid start \le pos < start + count,\ \operatorname{mod}(ci(pos), f) \neq 0 \,\big]
\end{aligned}
\end{equation*}

Ten verified lemmas in \texttt{GapProperties.scala} characterize this
sequence.

\subsection{Survivor Exactness}

$S$ is exact by construction: it is defined as exactly the non-multiples of
$f$ in range, nothing more and nothing less. The recursive Scala
implementation of \texttt{survivorValues} is verified to compute exactly
this specification.

{\raggedright
This is verified in
\href{https://github.com/thiagomata/prime-numbers/blob/master/src/main/scala/v1/chapter4/cycle/integral/recursive/properties/GapProperties.scala}{\texttt{GapProperties::\allowbreak assertSurvivorValuesContainsOnlyNonMultiples}}
and
\href{https://github.com/thiagomata/prime-numbers/blob/master/src/main/scala/v1/chapter4/cycle/integral/recursive/properties/GapProperties.scala}{\texttt{assertSurvivorValuesContainsNonMultipleAtPosition}}.
\par}

\subsection{Survivor Structure}

When every value in $[start, pos)$ is a multiple of $f$ and $ci(pos)$
itself is not, $ci(pos)$ is the first survivor, and $S$ decomposes as that
value followed by the survivors of the rest of the range, $(pos, end)$
where $end := start + count$:

\begin{equation*}
\begin{aligned}
&\big(\forall\, q \in [start, pos),\ \operatorname{mod}(ci(q), f) = 0\big)
  \;\land\; \operatorname{mod}(ci(pos), f) \neq 0 \\
&\quad \implies \operatorname{head}(S) = ci(pos)
  \quad \text{[First survivor]} \\
&S = ci(pos) :: \big[\, ci(q) \mid pos < q < end,\ \operatorname{mod}(ci(q), f) \neq 0 \,\big]
  \quad \text{[Structural split]}
\end{aligned}
\end{equation*}

At the endpoints of a range that itself survives, the first and last
elements of $S$ are exactly the endpoint values, and over one complete
period the difference between them equals the cycle's total gap sum ---
filtering does not change how far the integral advances across a full
period:

\begin{equation*}
\begin{aligned}
&\operatorname{mod}(ci(start), f) \neq 0
  \implies \operatorname{head}(S) = ci(start) \\
&\operatorname{mod}(ci(start + count - 1), f) \neq 0
  \implies \operatorname{last}(S) \\
&\quad = ci(start + count - 1) \\
&\operatorname{mod}(ci(0), f) \neq 0 \;\land\; \operatorname{mod}(ci(n), f) \neq 0
  \implies \operatorname{last}(S) - \operatorname{head}(S) = \operatorname{periodSum}(ci)
\end{aligned}
\end{equation*}

The gap between two consecutive survivors, once every intervening multiple
has been filtered out, is strictly positive:

\begin{equation*}
\begin{aligned}
&\operatorname{mod}(ci(from), f) \neq 0 \;\land\; \operatorname{mod}(ci(to), f) \neq 0
\;\land\; \big(\forall\, q \in (from, to),\ \operatorname{mod}(ci(q), f) = 0\big) \\
&\quad \implies ci(to) - ci(from) > 0
\end{aligned}
\end{equation*}

{\raggedright
All survivor-structure properties are verified in
\href{https://github.com/thiagomata/prime-numbers/blob/master/src/main/scala/v1/chapter4/cycle/integral/recursive/properties/GapProperties.scala}{\texttt{GapProperties::\allowbreak assertFirstSurvivorAtPosition}},
\href{https://github.com/thiagomata/prime-numbers/blob/master/src/main/scala/v1/chapter4/cycle/integral/recursive/properties/GapProperties.scala}{\texttt{assertSurvivorValuesSplitAtFirstPosition}},
\href{https://github.com/thiagomata/prime-numbers/blob/master/src/main/scala/v1/chapter4/cycle/integral/recursive/properties/GapProperties.scala}{\texttt{assertFirstSurvivorIsHead}},
\href{https://github.com/thiagomata/prime-numbers/blob/master/src/main/scala/v1/chapter4/cycle/integral/recursive/properties/GapProperties.scala}{\texttt{assertLastSurvivorIsLastScanned}},
\href{https://github.com/thiagomata/prime-numbers/blob/master/src/main/scala/v1/chapter4/cycle/integral/recursive/properties/GapProperties.scala}{\texttt{assertFilteredSumEqualsOriginalSum}},
and
\href{https://github.com/thiagomata/prime-numbers/blob/master/src/main/scala/v1/chapter4/cycle/integral/recursive/properties/GapProperties.scala}{\texttt{assertMergedGapPositive}}.
\par}

Filtering removes a value; the surviving values must still form a valid
cycle-integral gap list. This section proves that removing one multiple by
merging its two neighboring gaps produces exactly the cycle integral that
a fresh construction from the survivor list would produce, and that the
result contains no multiples of the filter value anywhere.

\subsection{Merge Shift Law}

Let $ci$ be a cycle integral with period $n$, and let $m$ be a merge index
with $0 \le m$ and $m + 1 < n$. Let $ci'$ be a cycle integral with period
$n - 1$, the same initial value, and a gap cycle equal to $ci$'s except
that the gaps at $m$ and $m + 1$ are combined into one:

\begin{equation*}
\begin{aligned}
\operatorname{cycle}'(i) &= \operatorname{cycle}(i)
  &&\text{for } i < m \\
\operatorname{cycle}'(m) &= \operatorname{cycle}(m) + \operatorname{cycle}(m + 1) \\
\operatorname{cycle}'(i) &= \operatorname{cycle}(i + 1)
  &&\text{for } i > m
\end{aligned}
\end{equation*}

Then $ci'$ agrees with $ci$ before the merge point, and agrees with $ci$
one position ahead at and after it:

\begin{equation*}
\begin{aligned}
pos < m &\implies ci'(pos) = ci(pos)
  &&\text{[Before merge]} \\
ci'(m) &= ci(m + 1)
  &&\text{[At merge]} \\
pos > m &\implies ci'(pos) = ci(pos + 1)
  &&\text{[After merge]} \\
ci'(n - 1) &= ci(n)
  &&\text{[Period boundary]}
\end{aligned}
\end{equation*}

\textbf{Proof.} By induction on the position, using the step property
(Subsection~3.1) in each of three cases split at the merge point.

\textbf{Case 1: Before the merge} ($pos < m$), by induction on $pos$.

\textbf{Base Case} ($pos = 0$):

\begin{equation*}
\begin{aligned}
ci'(0) &= init + \operatorname{cycle}'(0)
  &&\text{[By Definition]} \\
&= init + \operatorname{cycle}(0)
  &&\text{[Same initial value; unchanged gap before merge]} \\
&= ci(0)
  &&\text{[By Definition]}
\end{aligned}
\end{equation*}

\textbf{Induction Step} ($0 < pos < m$):

\begin{equation*}
\begin{aligned}
ci'(pos - 1) &= ci(pos - 1)
  &&\text{[Induction Hypothesis]} \\
ci'(pos) &= ci'(pos - 1) + \operatorname{cycle}'(pos)
  &&\text{[By Definition]} \\
&= ci(pos - 1) + \operatorname{cycle}(pos)
  &&\text{[Substitution; unchanged gap before merge]} \\
&= ci(pos)
  &&\text{[By Definition]}
\end{aligned}
\end{equation*}

\begin{equation*}
\therefore \ pos < m \implies ci'(pos) = ci(pos) \quad \blacksquare
\end{equation*}

\textbf{Case 2: At the merge} ($pos = m$). For $m > 0$, using Case 1 at
$m - 1$:

\begin{equation*}
\begin{aligned}
ci'(m) &= ci'(m - 1) + \operatorname{cycle}'(m)
  &&\text{[By Definition]} \\
&= ci(m - 1) + \big(\operatorname{cycle}(m) + \operatorname{cycle}(m + 1)\big)
  &&\text{[Case 1; merged-gap definition]} \\
&= \big(ci(m - 1) + \operatorname{cycle}(m)\big) + \operatorname{cycle}(m + 1)
  &&\text{[Regroup]} \\
&= ci(m) + \operatorname{cycle}(m + 1)
  &&\text{[By Definition]} \\
&= ci(m + 1)
  &&\text{[By Definition]}
\end{aligned}
\end{equation*}

For $m = 0$, the same identity holds directly, with no Case 1 dependency:
$ci'(0) = init + \operatorname{cycle}'(0) = init + \operatorname{cycle}(0) + \operatorname{cycle}(1) = ci(1)$.

\begin{equation*}
\therefore \ ci'(m) = ci(m + 1) \quad \blacksquare
\end{equation*}

\textbf{Case 3: After the merge} ($pos > m$), by induction on $pos$.

\textbf{Base Case} ($pos = m + 1$), using Case 2:

\begin{equation*}
\begin{aligned}
ci'(m + 1) &= ci'(m) + \operatorname{cycle}'(m + 1)
  &&\text{[By Definition]} \\
&= ci(m + 1) + \operatorname{cycle}(m + 2)
  &&\text{[Case 2; shifted-gap definition]} \\
&= ci(m + 2)
  &&\text{[By Definition]}
\end{aligned}
\end{equation*}

\textbf{Induction Step} ($pos > m + 1$):

\begin{equation*}
\begin{aligned}
ci'(pos - 1) &= ci(pos)
  &&\text{[Induction Hypothesis]} \\
ci'(pos) &= ci'(pos - 1) + \operatorname{cycle}'(pos)
  &&\text{[By Definition]} \\
&= ci(pos) + \operatorname{cycle}(pos + 1)
  &&\text{[Substitution; shifted-gap definition]} \\
&= ci(pos + 1)
  &&\text{[By Definition]}
\end{aligned}
\end{equation*}

\begin{equation*}
\therefore \ pos > m \implies ci'(pos) = ci(pos + 1) \quad \blacksquare
\end{equation*}

The period boundary case is the last instance of Case 3, at $pos = n - 1$:
$ci'(n - 1) = ci(n)$.

{\raggedright
This property is verified in
\href{https://github.com/thiagomata/prime-numbers/blob/master/src/main/scala/v1/chapter4/cycle/integral/recursive/properties/CycleIntegralFilterProperties.scala}{\texttt{CycleIntegralFilterProperties::\allowbreak assertSameBeforeMerge}},
\href{https://github.com/thiagomata/prime-numbers/blob/master/src/main/scala/v1/chapter4/cycle/integral/recursive/properties/CycleIntegralFilterProperties.scala}{\texttt{assertShiftAtMerge}},
\href{https://github.com/thiagomata/prime-numbers/blob/master/src/main/scala/v1/chapter4/cycle/integral/recursive/properties/CycleIntegralFilterProperties.scala}{\texttt{assertShiftAfterMerge}},
and
\href{https://github.com/thiagomata/prime-numbers/blob/master/src/main/scala/v1/chapter4/cycle/integral/recursive/properties/CycleIntegralFilterProperties.scala}{\texttt{assertNewCIAtSizeEqualsOld}}.
The full Scala verification code is in Appendix~A.13.
\par}

\subsection{Removing a Multiple}

When the value at the merge point is a multiple of a filter value $f$,
merging its two neighboring gaps removes it from the sequence entirely:
the new integral's value at that position is the old integral's next
value, which by construction is not itself a multiple.

\begin{equation*}
\begin{aligned}
\operatorname{mod}(ci(0), f) \neq 0 \;\land\; \operatorname{mod}(ci(m), f) = 0
&\implies ci'(m) = ci(m + 1)
&&\text{[Multiple removed]} \\
\operatorname{mod}(ci(m + 1), f) \neq 0
&\implies \operatorname{mod}(ci'(m), f) \neq 0
&&\text{[Result is not a multiple]}
\end{aligned}
\end{equation*}

\textbf{Proof.} The first identity is the merge shift law's "at merge" case
(Subsection~6.7) applied directly: $ci'(m) = ci(m + 1)$
regardless of whether $ci(m)$ was a multiple. The second follows because
equal values satisfy the same modulo condition.

{\raggedright
This property is verified in
\href{https://github.com/thiagomata/prime-numbers/blob/master/src/main/scala/v1/chapter4/cycle/integral/recursive/properties/CycleIntegralFilterProperties.scala}{\texttt{CycleIntegralFilterProperties::\allowbreak assertRemoveOneMultiple}}
and
\href{https://github.com/thiagomata/prime-numbers/blob/master/src/main/scala/v1/chapter4/cycle/integral/recursive/properties/CycleIntegralFilterProperties.scala}{\texttt{assertRemoveMultipleModNotZero}}.
The full Scala verification code is in Appendix~A.14.
\par}

\subsection{Direct Construction from Survivors}

Rather than merging gaps one multiple at a time, a filtered cycle integral
can be built directly from the survivor list (Subsection~6.5):
its initial value is the first survivor, and its gap cycle is the list of
differences between consecutive survivors. This direct construction
reproduces exactly what repeated merging would produce.

\begin{equation*}
\begin{aligned}
S &:= \operatorname{survivorValues}(ci, f, 0, n) \\
ci''\text{'s initial value} &= S_0,\quad \operatorname{cycle}''(i) = S_{i+1} - S_i \\
\implies ci''(k) &= S_{k + 1}
&&\text{[Direct construction matches survivors]}
\end{aligned}
\end{equation*}

\textbf{Proof.} By induction on $k$. The base case follows directly from the
integral's own definition applied to the first gap. The inductive step
adds the $k$-th constructed gap, which by definition equals
$S_{k+1} - S_k$, to the inductive hypothesis $ci''(k-1) = S_k$, giving
$ci''(k) = S_{k+1}$.

{\raggedright
This property is verified in
\href{https://github.com/thiagomata/prime-numbers/blob/master/src/main/scala/v1/chapter4/cycle/integral/recursive/properties/CycleIntegralFilterProperties.scala}{\texttt{CycleIntegralFilterProperties::\allowbreak assertNewCIGeneratesFiltered}},
\href{https://github.com/thiagomata/prime-numbers/blob/master/src/main/scala/v1/chapter4/cycle/integral/recursive/properties/CycleIntegralFilterProperties.scala}{\texttt{assertNewCIMatchesSurvivors}},
and
\href{https://github.com/thiagomata/prime-numbers/blob/master/src/main/scala/v1/chapter4/cycle/integral/recursive/properties/CycleIntegralFilterProperties.scala}{\texttt{assertGapsFromSurvivorsMatchCI}}.
The full Scala verification code is in Appendix~A.15.
\par}

\subsection{Filtered Result Has No Multiples}

The cycle integral constructed directly from the survivor list contains no
multiples of the filter value at any position --- its values are, by
construction, exactly the survivors.

\begin{equation*}
\operatorname{mod}(ci(0), f) \neq 0 \implies
\operatorname{mod}(ci''(k), f) \neq 0
\quad \text{for every valid } k
\end{equation*}

\textbf{Proof.} By Subsection~6.9, $ci''(k)$
equals the survivor $S_{k+1}$, and every element of the survivor list is
already known not to be a multiple of $f$
(Subsection~6.5).

\begin{equation*}
\begin{aligned}
ci''(k) &= S_{k+1} &&\text{[\S6.9]} \\
\operatorname{mod}(S_{k+1},f) &\neq 0 &&\text{[\S6.5]} \\
\therefore\ \operatorname{mod}(ci''(k),f) &\neq 0.
  \quad \blacksquare\ \text{[Q.E.D.]}
\end{aligned}
\end{equation*}

{\raggedright
This property is verified in
\href{https://github.com/thiagomata/prime-numbers/blob/master/src/main/scala/v1/chapter4/cycle/integral/recursive/properties/CycleIntegralFilterProperties.scala}{\texttt{CycleIntegralFilterProperties::\allowbreak assertFilterMergeComposition}}
and
\href{https://github.com/thiagomata/prime-numbers/blob/master/src/main/scala/v1/chapter4/cycle/integral/recursive/properties/CycleIntegralFilterProperties.scala}{\texttt{assertNextGapsValid}}.
The full Scala verification code is in Appendix~A.16.
\par}
